% ============================================================
%  preamble.tex —— 朴素中文笔记模板
%  编译方式：XeLaTeX，运行两遍以上以稳定目录与交叉引用
% ============================================================
\usepackage{geometry}
\geometry{a4paper, top=2.5cm, bottom=2.5cm, left=2.8cm, right=2.8cm}
\usepackage{amsmath,amssymb,amsthm}
\usepackage{graphicx}
\usepackage{float}
\usepackage{array}
\usepackage[hidelinks]{hyperref}

% 定理类环境：编号随章（如“定义 1.1”）
\newtheorem{theorem}{定理}[chapter]
\newtheorem{lemma}[theorem]{引理}
\newtheorem{corollary}[theorem]{推论}
\newtheorem{definition}[theorem]{定义}
\newtheorem{example}[theorem]{例题}
\renewcommand{\proofname}{证明}

% 代码排版（listings）：注意保持代码为 ASCII，中文说明写在代码环境之外
\usepackage{xcolor}
\usepackage{listings}

\definecolor{codebg}{RGB}{246,246,246}
\definecolor{codecomment}{RGB}{0,128,0}
\definecolor{codestring}{RGB}{163,21,21}
\definecolor{codekeyword}{RGB}{0,0,180}

\lstdefinestyle{code}{
	basicstyle=\ttfamily\small,
	backgroundcolor=\color{codebg},
	frame=single,
	framerule=0.4pt,
	rulecolor=\color{black!30},
	breaklines=true,
	columns=fullflexible,
	keepspaces=true,
	showstringspaces=false,
	tabsize=2,
	commentstyle=\color{codecomment}\itshape,
	stringstyle=\color{codestring},
	keywordstyle=\color{codekeyword}
}

\lstdefinestyle{xml}{
	style=code,
	language=XML,
	morecomment=[s]{<!--}{-->}
}
